\hypertarget{thiux1ebft-lux1eadp-phuxe1t-triux1ec3n-cux1ee5c-bux1ed9}{%
\section{Thiết lập phát triển cục
bộ}}

Phần này bao gồm việc thiết lập môi trường phát triển cục bộ hoàn chỉnh
sử dụng Docker Compose để mô phỏng các dịch vụ AWS cục bộ.

\hypertarget{tux1ed5ng-quan-kiux1ebfn-truxfac-local}{%
\subsection{Tổng quan kiến trúc
(Local)}}

\emph{(Mã nguồn chi tiết đã được lược bỏ để báo cáo ngắn gọn. Vui lòng
xem trong source code dự án)}

\hypertarget{cuxe1c-dux1ecbch-vux1ee5-docker-compose}{%
\subsection{Các dịch vụ Docker
Compose}}

File \texttt{docker-compose.yml} định nghĩa tất cả các dịch vụ cần thiết
cho phát triển cục bộ:

\emph{(Mã nguồn chi tiết đã được lược bỏ để báo cáo ngắn gọn. Vui lòng
xem trong source code dự án)}

\hypertarget{khux1edfi-ux111ux1ed9ng-muxf4i-trux1b0ux1eddng}{%
\subsection{Khởi động môi
trường}}

\begin{Shaded}
\begin{Highlighting}[]
\CommentTok{\# Khởi động tất cả dịch vụ nền}
\ExtensionTok{docker}\NormalTok{ compose up }\AttributeTok{{-}d}

\CommentTok{\# Kiểm tra trạng thái}
\ExtensionTok{docker}\NormalTok{ compose ps}

\CommentTok{\# Xem logs}
\ExtensionTok{docker}\NormalTok{ compose logs }\AttributeTok{{-}f}\NormalTok{ postgres}
\ExtensionTok{docker}\NormalTok{ compose logs }\AttributeTok{{-}f}\NormalTok{ kafka}
\end{Highlighting}
\end{Shaded}

Kết quả mong đợi:

\begin{verbatim}
NAME                 STATUS
newsrag-postgres     Up (healthy)
newsrag-qdrant       Up
newsrag-zookeeper    Up
newsrag-kafka        Up (healthy)
newsrag-kafka-ui     Up
\end{verbatim}

\hypertarget{xuxe1c-minh-cuxe1c-dux1ecbch-vux1ee5}{%
\subsection{Xác minh các dịch
vụ}}

\hypertarget{postgresql}{%
\subsubsection{PostgreSQL}}

\begin{Shaded}
\begin{Highlighting}[]
\CommentTok{\# Kết nối database}
\ExtensionTok{psql} \AttributeTok{{-}h}\NormalTok{ localhost }\AttributeTok{{-}U}\NormalTok{ postgres }\AttributeTok{{-}d}\NormalTok{ newsrag}

\CommentTok{\# Hoặc qua Docker}
\ExtensionTok{docker}\NormalTok{ exec }\AttributeTok{{-}it}\NormalTok{ newsrag{-}postgres psql }\AttributeTok{{-}U}\NormalTok{ postgres }\AttributeTok{{-}d}\NormalTok{ newsrag}

\CommentTok{\# Xác minh tables được tạo từ warehouse.sql}
\ExtensionTok{\textbackslash{}dt}
\end{Highlighting}
\end{Shaded}

Tables mong đợi:

\begin{verbatim}
         List of relations
 Schema |       Name       | Type  |  Owner
--------+------------------+-------+----------
 public | dim_author       | table | postgres
 public | dim_content      | table | postgres
 public | dim_source       | table | postgres
 public | dim_time         | table | postgres
 public | fact_article_authors | table | postgres
 public | fact_articles    | table | postgres
 public | fact_chunks      | table | postgres
\end{verbatim}

\hypertarget{qdrant}{%
\subsubsection{Qdrant}}

\begin{Shaded}
\begin{Highlighting}[]
\CommentTok{\# Kiểm tra collections}
\ExtensionTok{curl}\NormalTok{ http://localhost:6333/collections}

\CommentTok{\# Kết quả mong đợi: \{"result":\{"collections":[]\},"status":"ok","time":0.0...\}}
\end{Highlighting}
\end{Shaded}

\hypertarget{kafka}{%
\subsubsection{Kafka}}

\begin{Shaded}
\begin{Highlighting}[]
\CommentTok{\# Liệt kê topics}
\ExtensionTok{docker}\NormalTok{ exec newsrag{-}kafka kafka{-}topics }\AttributeTok{{-}{-}bootstrap{-}server}\NormalTok{ localhost:9092 }\AttributeTok{{-}{-}list}

\CommentTok{\# Tạo news\_raw topic (nếu chưa auto{-}created)}
\ExtensionTok{docker}\NormalTok{ exec newsrag{-}kafka kafka{-}topics }\AttributeTok{{-}{-}bootstrap{-}server}\NormalTok{ localhost:9092 }\AttributeTok{{-}{-}create} \AttributeTok{{-}{-}topic}\NormalTok{ news\_raw }\AttributeTok{{-}{-}partitions}\NormalTok{ 3 }\AttributeTok{{-}{-}replication{-}factor}\NormalTok{ 1}
\end{Highlighting}
\end{Shaded}

\hypertarget{kafka-ui}{%
\subsubsection{Kafka UI}}

Mở http://localhost:8080 trong trình duyệt để xem topics, messages,
consumer groups.

\hypertarget{thiux1ebft-lux1eadp-muxf4i-trux1b0ux1eddng-python}{%
\subsection{Thiết lập môi trường
Python}}

\emph{(Mã nguồn chi tiết đã được lược bỏ để báo cáo ngắn gọn. Vui lòng
xem trong source code dự án)}

\hypertarget{requirements.txt}{%
\subsubsection{requirements.txt}}

\emph{(Mã nguồn chi tiết đã được lược bỏ để báo cáo ngắn gọn. Vui lòng
xem trong source code dự án)}

\hypertarget{cux1ea5u-huxecnh-biux1ebfn-muxf4i-trux1b0ux1eddng}{%
\subsection{Cấu hình biến môi
trường}}

\emph{(Mã nguồn chi tiết đã được lược bỏ để báo cáo ngắn gọn. Vui lòng
xem trong source code dự án)}

\hypertarget{chux1ea1y-pipeline-cux1ee5c-bux1ed9}{%
\subsection{Chạy Pipeline cục
bộ}}

\hypertarget{crawl-tin-tux1ee9c-scrapy-rightarrow-kafka}{%
\subsubsection{\texorpdfstring{1. Crawl tin tức (Scrapy \(\rightarrow\)
Kafka)}{1. Crawl tin tức (Scrapy \textbackslash rightarrow Kafka)}}}

\begin{Shaded}
\begin{Highlighting}[]
\CommentTok{\# Terminal 1: Khởi động Kafka Consumer (ghi vào PostgreSQL)}
\ExtensionTok{python}\NormalTok{ consumer/consumer.py}

\CommentTok{\# Terminal 2: Chạy Crawler}
\ExtensionTok{python}\NormalTok{ main.py }\AttributeTok{{-}{-}mode}\NormalTok{ crawl}
\end{Highlighting}
\end{Shaded}

Crawler: - Đọc \texttt{config/config\_site.json} để lấy danh sách URL
báo - Chạy Scrapy SitemapSpider cho mỗi trang - Đẩy dữ liệu bài viết vào
Kafka topic \texttt{news\_raw}

\hypertarget{chux1ea1y-etl-postgresql-rightarrow-star-schema}{%
\subsubsection{\texorpdfstring{2. Chạy ETL (PostgreSQL \(\rightarrow\)
Star
Schema)}{2. Chạy ETL (PostgreSQL \textbackslash rightarrow Star Schema)}}}

\begin{Shaded}
\begin{Highlighting}[]
\CommentTok{\# Sau khi crawler xong và consumer đã xử lý messages}
\ExtensionTok{python}\NormalTok{ main.py }\AttributeTok{{-}{-}mode}\NormalTok{ etl}
\end{Highlighting}
\end{Shaded}

ETL process: - Đọc bài thô từ bảng \texttt{article\_metadata} - Làm sạch
HTML, trích xuất text - Chunk text thành các đoạn \textasciitilde500
token - Điền Star Schema tables (\texttt{dim\_*}, \texttt{fact\_*}) -
Cập nhật \texttt{embedded\ =\ true} trên bài đã xử lý

\hypertarget{vectorize-star-schema-rightarrow-qdrant}{%
\subsubsection{\texorpdfstring{3. Vectorize (Star Schema \(\rightarrow\)
Qdrant)}{3. Vectorize (Star Schema \textbackslash rightarrow Qdrant)}}}

\begin{Shaded}
\begin{Highlighting}[]
\ExtensionTok{python}\NormalTok{ main.py }\AttributeTok{{-}{-}mode}\NormalTok{ vectorize}
\end{Highlighting}
\end{Shaded}

Vectorizer: - Đọc chunks từ bảng \texttt{fact\_chunks} - Tạo embeddings
bằng SentenceTransformer (model local) - Upsert vectors vào Qdrant
collection \texttt{news\_chunks}

\hypertarget{chux1ea1y-full-pipeline-tux1ef1-ux111ux1ed9ng}{%
\subsubsection{4. Chạy Full Pipeline (Tự
động)}}

\begin{Shaded}
\begin{Highlighting}[]
\CommentTok{\# Chạy crawl $\textbackslash{}rightarrow$ ETL $\textbackslash{}rightarrow$ vectorize lần lượt}
\ExtensionTok{python}\NormalTok{ main.py }\AttributeTok{{-}{-}mode}\NormalTok{ full}
\end{Highlighting}
\end{Shaded}

\hypertarget{test-rag-api-interactive}{%
\subsubsection{5. Test RAG API
(Interactive)}}

\begin{Shaded}
\begin{Highlighting}[]
\CommentTok{\# Khởi động FastAPI server}
\ExtensionTok{uvicorn}\NormalTok{ search.api:app }\AttributeTok{{-}{-}reload} \AttributeTok{{-}{-}port}\NormalTok{ 8000}

\CommentTok{\# Hoặc dùng Makefile}
\FunctionTok{make}\NormalTok{ test{-}interactive}
\end{Highlighting}
\end{Shaded}

Test trên trình duyệt: http://localhost:8000/docs (Swagger UI)

Hoặc CLI:

\begin{Shaded}
\begin{Highlighting}[]
\ExtensionTok{python} \AttributeTok{{-}c} \StringTok{"}
\StringTok{from search.engine import Pipeline}
\StringTok{p = Pipeline()}
\StringTok{result = p.ask(\textquotesingle{}Tóm tắt tin tức về kinh tế Việt Nam hôm nay\textquotesingle{})}
\StringTok{print(result.summary)}
\StringTok{print(f\textquotesingle{}Sources: \{result.total\}, Time: \{result.duration\_ms\}ms\textquotesingle{})}
\StringTok{"}
\end{Highlighting}
\end{Shaded}

\hypertarget{makefile-commands}{%
\subsection{Makefile Commands}}

\emph{(Mã nguồn chi tiết đã được lược bỏ để báo cáo ngắn gọn. Vui lòng
xem trong source code dự án)}

\hypertarget{xux1eed-luxfd-sux1ef1-cux1ed1}{%
\subsection{Xử lý sự cố}}

\begingroup
\small
\setlength{\tabcolsep}{3pt}
\renewcommand{\arraystretch}{1.15}
\begin{longtable}[]{@{}
  >{\raggedright\arraybackslash}p{(\columnwidth - 2\tabcolsep) * \real{0.4211}}
  >{\raggedright\arraybackslash}p{(\columnwidth - 2\tabcolsep) * \real{0.5789}}@{}}
\toprule\noalign{}
\begin{minipage}[b]{\linewidth}\raggedright
Vấn đề
\end{minipage} & \begin{minipage}[b]{\linewidth}\raggedright
Giải pháp
\end{minipage} \\
\midrule\noalign{}
\endhead
\bottomrule\noalign{}
\endlastfoot
\texttt{PostgreSQL\ connection\ refused} & Kiểm tra
\texttt{docker\ compose\ ps\ postgres}, chờ healthy. Xem
\texttt{docker\ compose\ logs\ postgres} \\
\texttt{Kafka\ not\ ready} & Đợi Kafka healthy (30-60s). Kiểm tra
\texttt{docker\ compose\ logs\ -f\ kafka} \\
\texttt{Qdrant\ collection\ issues} & Xóa và tạo lại collection:
\texttt{curl\ -X\ DELETE\ http://localhost:6333/collections/news\_chunks} \\
\texttt{SentenceTransformer\ model\ download} & Lần đầu tải model
(\textasciitilde100MB). Lần sau dùng cache tại
\texttt{\textasciitilde{}/.cache/huggingface/} \\
\texttt{Port\ conflicts} & Kiểm tra port 5432, 6333, 9092, 8080. Dừng
dịch vụ xung đột hoặc đổi port trong docker-compose.yml \\
\end{longtable}
\endgroup

\hypertarget{quy-truxecnh-phuxe1t-triux1ec3n}{%
\subsection{Quy trình phát
triển}}

\begin{verbatim}
1. Sửa code (crawler/spiders/spider.py, etl/etl_warehouse.py, etc.)
2. Test cục bộ: make crawl / make etl / make vectorize
3. Kiểm tra logs: docker compose logs -f <service>
4. Xác minh DB: psql -h localhost -U postgres -d newsrag -c "SELECT * FROM fact_chunks LIMIT 5;"
5. Test RAG: make test-interactive
6. Commit changes
7. Build Docker image cho AWS deployment (xem 5.9-AWS-Deploy)
\end{verbatim}

\hypertarget{bux1b0ux1edbc-tiux1ebfp-theo}{%
\subsection{Bước tiếp theo}}

Sau khi phát triển cục bộ hoạt động: 1. \textbf{Build Docker image} cho
AWS: \href{5.9-AWS-Deploy/}{AWS Deployment Prep} 2. \textbf{Deploy lên
AWS} --- \href{5.10-Fargate-Crawler/}{Fargate Crawler},
\href{5.11-Lambda-Consumer/}{Lambda Consumer},
\href{5.12-Lambda-ETL/}{Lambda ETL} 3. \textbf{Test RAG API} ---
\href{5.13-RAG-API/}{RAG API}

\begin{center}\rule{0.5\linewidth}{0.5pt}\end{center}

\textbf{Tiếp theo:} \href{5.5-Crawler/}{Phát triển Crawler}
